\chapter{Hand Pain}

\section*{Definition}
Hand pain refers to discomfort, soreness, or aching in any part of the hand, including the fingers, knuckles, palm, and wrist area. It can range from a dull ache to sharp, stabbing pain and may be caused by injury, overuse, or underlying medical conditions.

\section*{What Happens in Your Body}
Your hand contains 27 bones, numerous joints, ligaments, tendons, nerves, and blood vessels — all working together for fine motor function. Pain can arise from any of these structures. Inflammation of the joints (arthritis) causes stiffness and aching. Tendon irritation (tendonitis) produces pain with movement. Nerve compression (such as carpal tunnel syndrome) causes tingling, numbness, or burning pain. Injury to bones or ligaments causes localized sharp pain and swelling.

\section*{Causes}
\begin{itemize}
  \item Overuse and repetitive strain: Prolonged typing, hand-intensive work, or repetitive gripping
  \item Arthritis: Osteoarthritis (wear-and-tear), rheumatoid arthritis (autoimmune), psoriatic arthritis, gout
  \item Carpal tunnel syndrome: Compression of the median nerve in the wrist, causing pain, numbness, and tingling in the thumb and first three fingers
  \item De Quervain's tenosynovitis: Pain on the thumb side of the wrist from inflamed tendons
  \item Trigger finger: A finger that catches or locks in a bent position due to tendon sheath inflammation
  \item Hand fractures or sprains: From falls or direct impact
  \item Peripheral neuropathy: From diabetes, vitamin B deficiency, or alcohol use
  \item Ganglion cysts: Fluid-filled sacs near joints or tendons that can cause pressure and pain
  \item Raynaud's phenomenon: Blood vessel spasms that cause pain, coldness, and color changes in fingers
  \item Infection: Cellulitis, paronychia (nail fold infection), or felon (finger pulp infection)
  \item Dupuytren's contracture: Thickening of tissue in the palm that can cause finger bending and discomfort
\end{itemize}

\section*{When to See a Doctor}
See your doctor if:
\begin{itemize}
  \item Hand pain follows an injury, especially if swelling is severe or you cannot move your fingers
  \item Pain is severe, persistent, or getting worse over time
  \item You have numbness, tingling, or weakness in your hand that does not go away
  \item Your fingers look pale, blue, or feel cold (possible blood flow problem)
  \item You have signs of infection: Redness, warmth, swelling, or pus
  \item The pain interferes with daily activities such as writing, dressing, or eating
  \item You have a known medical condition like diabetes or rheumatoid arthritis and develop new hand pain
\end{itemize}

\section*{Related Symptoms}
\begin{itemize}
  \item Swelling in the fingers or knuckles
  \item Stiffness, especially in the morning
  \item Numbness or tingling (pins and needles)
  \item Weakness or dropping objects
  \item Clicking or popping sensation with finger movement
  \item Changes in skin color (pale, blue, red)
  \item Visible lumps or bumps on the hand or wrist
\end{itemize}
\textbf{Important:} Sudden hand pain with chest pressure, arm pain, or shortness of breath could be a sign of a heart attack — call emergency services.

\section*{Self-Care / Home Management}
\begin{itemize}
  \item Rest the affected hand and avoid activities that trigger the pain.
  \item Apply ice for 15-20 minutes several times a day if there is swelling.
  \item Apply heat for stiffness without swelling.
  \item Take over-the-counter pain relievers such as ibuprofen (Advil, Motrin) or acetaminophen (Tylenol).
  \item Wear a splint or brace if recommended, especially at night for carpal tunnel symptoms.
  \item Perform gentle stretching exercises if approved by your doctor.
\end{itemize}

\section*{How Your Doctor Will Evaluate You}
\begin{itemize}
  \item Your doctor will ask about your occupation, hobbies, injury history, and when the pain occurs.
  \item They will examine your hand for swelling, tenderness, range of motion, and strength.
  \item They may perform specific tests such as Tinel's sign (tapping over the wrist for carpal tunnel) or Finkelstein's test (for De Quervain's).
  \item X-rays can show fractures, arthritis, or bone abnormalities.
  \item MRI or ultrasound may be used to evaluate tendons, ligaments, and soft tissue.
  \item Nerve conduction studies can confirm carpal tunnel syndrome or other nerve compression.
  \item Blood tests may be ordered if rheumatoid arthritis, gout, or infection is suspected.
\end{itemize}

\section*{Treatment Options}
\begin{itemize}
  \item Rest, ice, compression, and elevation (RICE) for acute injuries.
  \item Anti-inflammatory medications (NSAIDs like ibuprofen or naproxen) for arthritis and tendonitis.
  \item Splints or braces to rest the affected structures — especially for carpal tunnel syndrome and trigger finger.
  \item Corticosteroid injections for inflamed joints or tendons.
  \item Physical therapy or hand therapy exercises to restore strength and range of motion.
  \item Surgery for severe carpal tunnel syndrome, trigger finger that does not respond to treatment, or Dupuytren's contracture.
  \item Disease-modifying medications for rheumatoid arthritis or psoriatic arthritis.
  \item Antibiotics for infections.
\end{itemize}

\section*{References}

\begin{thebibliography}{99}
\bibitem{harrison-msk} Longo DL, et al. \emph{Harrison's Principles of Internal Medicine}. 21st ed. McGraw-Hill; 2022. Chapter 390: Approach to the Patient with Musculoskeletal Pain.
\bibitem{uptodate-hand} Shmerling RH, et al. Evaluation of the patient with hand pain. In: UpToDate, Post TW (Ed), UpToDate, Waltham, MA; 2025.
\bibitem{wolfe} Wolfe F, et al. Classification of hand osteoarthritis. \emph{Arthritis Care Res}. 2022;74(8):1219-1230.
\bibitem{green} Green DP, et al. \emph{Green's Operative Hand Surgery}. 8th ed. Elsevier; 2022. Chapter 1: Examination of the Hand.
\bibitem{aaos-hand} American Academy of Orthopaedic Surgeons. \emph{AAOS Clinical Practice Guideline: Treatment of Osteoarthritis of the Hand}. AAOS; 2023.
\bibitem{macdermid} MacDermid JC, et al. Clinical assessment of hand pain and function. \emph{J Hand Ther}. 2023;36(2):295-308.
\end{thebibliography}
